%==============================================================================
% Proposta de Trabalho de Conclusao de Curso
% Modelo aberto de conversor HVDC-VSC (MMC meia-ponte) para ATP-EMTP
%
% Compila em pdfLaTeX no Overleaf, sem arquivos auxiliares.
% Menu do Overleaf: Compiler = pdfLaTeX.
%==============================================================================
\documentclass[12pt,a4paper]{article}

%------------------------------------------------------------------ idioma ---
\usepackage[T1]{fontenc}
\usepackage[utf8]{inputenc}
\usepackage[brazil]{babel}

%------------------------------------------------------------------ fontes ---
% Times-like (padrao mais aceito em normas ABNT).
% Para voltar a fonte padrao do LaTeX, comente a linha abaixo.
\usepackage{newtxtext,newtxmath}
\usepackage{microtype}

%------------------------------------------------------------------ layout ---
\usepackage[left=3cm,right=2cm,top=3cm,bottom=2cm]{geometry}
\usepackage{setspace}
\usepackage{indentfirst}
\usepackage{ragged2e}

%------------------------------------------------------------------ tabelas --
\usepackage{booktabs}
\usepackage{tabularx}
\usepackage{longtable}
\usepackage{xltabular}
\usepackage{array}
\usepackage{multirow}

%------------------------------------------------------------------- outros --
\usepackage{xcolor}
\usepackage{enumitem}
\usepackage{titlesec}
\usepackage{fancyhdr}
\usepackage{csquotes}
\usepackage{textcomp}
\usepackage{hyperref}

%--------------------------------------------------------------- aparencia ---
\definecolor{tinta}{HTML}{1A1A1A}
\definecolor{acento}{HTML}{28418C}
\definecolor{cobre}{HTML}{9C4A1E}
\definecolor{regra}{HTML}{BFC6D2}

\hypersetup{
  colorlinks=true,
  linkcolor=acento,
  urlcolor=acento,
  citecolor=acento,
  pdftitle={Proposta de TCC -- Modelo aberto de conversor HVDC-VSC para ATP-EMTP},
  pdfsubject={Modelagem EMT de conversor VSC-MMC meia-ponte em ATP-EMTP},
  pdfkeywords={HVDC-VSC, MMC, ATP-EMTP, transitorios eletromagneticos, EPE, ONS}
}

\titleformat{\section}{\normalfont\large\bfseries\color{acento}}{\thesection}{0.8em}{}
\titleformat{\subsection}{\normalfont\normalsize\bfseries}{\thesubsection}{0.7em}{}
\titlespacing*{\section}{0pt}{2.2ex plus 1ex minus .2ex}{1.4ex plus .2ex}

\pagestyle{fancy}
\fancyhf{}
\fancyhead[L]{\footnotesize Proposta de TCC \textbar{} Modelo aberto HVDC-VSC para ATP-EMTP}
\fancyhead[R]{\footnotesize\thepage}
\renewcommand{\headrulewidth}{0.4pt}

\setlength{\parindent}{1.25cm}
\onehalfspacing

% listas mais compactas
\setlist[itemize]{itemsep=2pt,topsep=4pt,parsep=0pt}
\setlist[enumerate]{itemsep=2pt,topsep=4pt,parsep=0pt}

% citacao longa no estilo ABNT (recuo 4 cm, corpo menor, espaco simples)
\newenvironment{citacao}
  {\begin{quote}\singlespacing\small\setlength{\leftskip}{2.5cm}\noindent}
  {\par\end{quote}}

% fonte da citacao
\newcommand{\fonte}[1]{\par\vspace{2pt}\noindent{\footnotesize\textcolor{gray}{#1}}}

% caixa de destaque simples (sem dependencia de tcolorbox)
\newenvironment{destaque}[1]
  {\par\vspace{6pt}\begin{list}{}{%
     \setlength{\leftmargin}{0.8cm}\setlength{\rightmargin}{0cm}%
     \setlength{\labelwidth}{0pt}\setlength{\labelsep}{0pt}}%
   \item[]\textcolor{acento}{\rule[-1pt]{2pt}{10pt}}\ \textbf{#1}\par\vspace{2pt}\small}
  {\end{list}\vspace{6pt}}

% caminho de arquivo / identificador de codigo
\newcommand{\arq}[1]{\texttt{\small #1}}

\newcommand{\pag}[1]{{\footnotesize\textcolor{gray}{p.~#1}}}

\renewcommand{\arraystretch}{1.22}

%==============================================================================
\begin{document}
%==============================================================================

%------------------------------------------------------------------- capa ----
\begin{titlepage}
\centering
\vspace*{1.5cm}

{\footnotesize\textsc{Proposta de Trabalho de Conclusão de Curso}\par}
\vspace{0.4cm}
{\footnotesize Engenharia Elétrica\par}

\vspace{2.5cm}

{\LARGE\bfseries Modelo aberto de conversor HVDC-VSC\\[0.25cm]
(MMC meia-ponte) para ATP-EMTP\par}

\vspace{0.9cm}

{\large Benchmark reprodutível do bipolo $\pm$600~kV Angicos--Itaporanga~2\\
e integração ao Olivas Power System Studio\par}

\vspace{2.5cm}

\begin{minipage}{0.92\textwidth}
\small
\renewcommand{\arraystretch}{1.35}
\begin{tabularx}{\textwidth}{@{}>{\bfseries}p{4.2cm}X@{}}
\toprule
Documento-base   & EPE-DEE-RE-071/2025-rev0 --- \textit{Estudo de Expansão das Interligações Regionais, Parte III} (nov.~2025, 509~p.) \\
Caso de estudo   & Bipolo Nordeste~II --- $\pm$600~kV, 3.000~MW, $\sim$2.500~km \\
Ferramenta       & ATP-EMTP + linguagem MODELS, integrado ao Olivas EPS \\
Empresa parceira & Olivas --- produto Olivas Power System Studio \\
Período previsto & set./2026 a jun./2027 (TCC~1 e TCC~2) \\
Autor            & \rule{5cm}{0.4pt} \\
Orientador(a)    & \rule{5cm}{0.4pt} \\
\bottomrule
\end{tabularx}
\end{minipage}

\vfill
{\footnotesize 27 de agosto de 2026\par}
\end{titlepage}

%------------------------------------------------------------------ sumario --
\pagenumbering{roman}
\tableofcontents
\clearpage
\pagenumbering{arabic}

%==============================================================================
\section{Sumário executivo}
%==============================================================================

O relatório R1 da Empresa de Pesquisa Energética recomenda o \textbf{bipolo Nordeste~II}:
HVDC-VSC, $\pm$600~kV, 3.000~MW, cerca de 2.500~km de linha aérea entre Angicos (RN) e
Itaporanga~2 (SP), com entrada em operação em 2033, dentro de um pacote de
\textbf{R\$~26,5 bilhões} --- dos quais R\$~17,1 bilhões, ou 65\%, em instalações de corrente
contínua. Será o primeiro elo VSC do Sistema Interligado Nacional; os seis bipolos em operação e o
sétimo em implantação são todos LCC.

No mesmo documento, a EPE registra por escrito a lacuna que dá origem a este trabalho:

\begin{citacao}
\enquote{A ausência de modelos oficiais validados para conversores VSC constituiu um desafio
relevante à realização das análises dinâmicas deste estudo. Tal fato dificultou, inclusive, o
desenvolvimento de \textbf{modelos genéricos}\footnote{Nota de rodapé 8 do relatório:
\enquote{Modelo concebido com os elementos contidos nas concepções atualizadas de VSC dos
principais fabricantes, de forma a não privilegiar concepções específicas.}}, em virtude da
\textbf{ausência de um benchmark} a ser utilizado como referência para comparação de resultados.}
\fonte{EPE-DEE-RE-071/2025-rev0, p.~21, \S1.5.5.}
\end{citacao}

A nota de rodapé 8 do relatório é, literalmente, a especificação do produto proposto aqui.

\subsection{Tema recomendado}

Desenvolver, validar e publicar sob licença permissiva um modelo genérico de estação conversora
HVDC-VSC com submódulos 100\% meia-ponte, executável em \textbf{ATP-EMTP}, acompanhado de um
\textbf{benchmark aberto e reprodutível} do bipolo Nordeste~II, e integrá-lo ao Olivas EPS como
componente de catálogo, emissor de cartões, estudo automatizado e laudo auditável.

\subsection{Resposta aos dois pontos levantados}

\begin{destaque}{Ponto 1 --- modelo público integrável ao Olivas EPS}
\textbf{Sim, com três precisões.} O termo correto é \emph{aberto sob licença permissiva}
(Apache-2.0 e CC-BY-4.0), e não \enquote{domínio público}: a Lei 9.610/98 não admite renúncia a
direitos morais, e instrumentos de dedicação ao domínio público são de tradição estrangeira. O
entregável é um \textbf{par}: um artefato que roda em ATP puro, independente do produto, e uma
camada de integração no Olivas EPS. E o gargalo real não é o MMC: é o acoplamento automático entre
MODELS e a rede, que hoje \textbf{não existe} no produto (Seção~\ref{sec:integracao}). Ele é a
primeira tarefa técnica e o gate de viabilidade.
\end{destaque}

\begin{destaque}{Ponto 2 --- viabilidade de mercado}
\textbf{Sim, por um caminho diferente do imaginado.} A premissa \enquote{só o PSCAD tem modelo
VSC} é \textbf{falsa} como afirmação geral e \textbf{verdadeira} em três recortes. O recorte
decisivo foi verificado na fonte primária: o ONS \textbf{obriga a entrega de modelo detalhado do
elo CC em ATP}, mesmo quando o estudo foi feito em PSCAD, e determina que os estudos de manobra do
pátio CA das conversoras sejam realizados \textbf{na ferramenta ATP}. Existe obrigação regulatória
de produzir modelo HVDC em ATP e não existe modelo VSC público em ATP. O dimensionamento honesto do
mercado está na Seção~\ref{sec:mercado}.
\end{destaque}

%==============================================================================
\section{Contexto: o que a EPE decidiu}
%==============================================================================

\subsection{A solução de referência}

\begin{center}
\small
\begin{tabularx}{\textwidth}{@{}>{\raggedright\arraybackslash}p{4.1cm}Xr@{}}
\toprule
\textbf{Item} & \textbf{Valor} & \textbf{Pág.} \\
\midrule
Denominação            & Bipolo Nordeste~II --- LT 600~kV CC Angicos (RN) -- Itaporanga~2 (SP) & 25 \\
Tecnologia             & HVDC-VSC, MMC \textbf{100\% Half-Bridge}                              & 20, 251 \\
Tensão e potência      & $\pm$600~kV, 3.000~MW (inversora 2.880~MW; reversa 2.800~MW)          & 84, 251 \\
Extensão               & $\sim$2.500~km em linha aérea                                          & 25, 252 \\
Condutor               & 6 $\times$ CAA 2167~MCM (Kiwi) por polo; 17~m entre polos; 14~m ao solo & 250, 252 \\
Resistência do polo    & 0,00425~$\Omega$/km a 20~\textcelsius; 0,00476~$\Omega$/km a 50~\textcelsius & 252 \\
Corrente e perdas na LT & 2,5~kA direta; 2,34~kA reversa; 132,7~MW                              & 252 \\
Controle               & \emph{Grid-Following} nas simulações; \emph{Grid-Forming} como requisito & 118, 42 \\
Modos de operação      & bipolar; monopolar com retorno por eletrodo; monopolar com retorno metálico & 42 \\
SCR de projeto         & operação exigida entre \textbf{1,5 e 6,5}                              & 43 \\
SCR calculado          & Angicos 4,08--6,07 (\textbf{1,92} na N-2 Açu~III--Angicos); Itaporanga~2 5,47--5,64 & 263--265 \\
SCR do modelo EMT      & equivalente de Thévenin com SCR = 2,5                                  & 205 \\
Investimento           & R\$~26,5 bi --- R\$~9,7 bi na LT CC e R\$~3,7 bi por conversora        & 23 \\
Entrada em operação    & 2033 (obras determinativas)                                            & 25 \\
\bottomrule
\end{tabularx}
\end{center}

Requisitos funcionais inéditos para o SIN (\S3.4[8], p.~42--43): modo \textbf{STATCOM}, injeção de
corrente reativa sob defeito, \textbf{grid-forming}, \textbf{black-start bidirecional}, transmissão
em potência reversa, operação a 85\% da tensão nominal e sobrecarga de 400~MW pela capacidade
inerente dos submódulos.

\subsection{Requisitos de desempenho que ainda não existem em norma}

O item \S3.4[9] exige reenergização do bipolo em \textbf{$\leq$~2,5~s} após a detecção da falta CC e
recuperação de 90\% da potência em \textbf{$\leq$~5~s}, e reconhece:

\begin{citacao}
\enquote{será necessária a flexibilização dos tempos atualmente previstos em procedimentos de rede
para elos de corrente contínua (que abarcam apenas a tecnologia LCC) --- que são da ordem de
350~ms --- ou a criação de procedimentos específicos para instalações VSC.}
\fonte{EPE-DEE-RE-071/2025-rev0, p.~43.}
\end{citacao}

Um requisito normativo novo terá de ser escrito, e sua verificação é, por natureza, uma simulação
de transitórios eletromagnéticos.

%==============================================================================
\section{O problema}
%==============================================================================

\subsection{O encadeamento de contorno adotado pela EPE}

\begin{center}
\small
\begin{tabularx}{\textwidth}{@{}p{3.1cm}p{4.4cm}X@{}}
\toprule
\textbf{Etapa} & \textbf{Representação usada} & \textbf{Limitação declarada} \\
\midrule
Fluxo de potência &
Barras PV com limites de reativo de $\sim$30\% da potência nominal &
Simplificação (p.~118--119) \\
Estabilidade eletromecânica (RMS) &
\textbf{Carga ZIP: 60\% de potência constante e 40\% de corrente constante} no ANATEM &
\enquote{não se mostrou suficiente para a realização de análises elétricas conclusivas nem para a
recomendação de reforços sistêmicos} (p.~119) \\
Estabilidade, segunda geração &
Modelo de fasores dinâmicos do MMC, \textbf{contratado ao CEPEL} &
\enquote{ainda não estivesse totalmente apto} a faltas CC (p.~119) \\
\textbf{Transitórios (EMT)} &
\textbf{PSCAD/EMTDC}, a partir do benchmark \emph{VSC-MMC Half Bridge} de 320~kV e 1.200~MW,
reescalado para 600~kV e 3.000~MW com apoio dos fabricantes &
Rede CA representada apenas por equivalentes de Thévenin (p.~206) \\
\bottomrule
\end{tabularx}
\end{center}

O desfecho, na conclusão do Capítulo 12 do relatório:

\begin{citacao}
\enquote{Atualmente, o modelo desenvolvido em ANATEM para representação do bipolo não representa de
forma adequada este tipo de dinâmica, portanto não foi possível avaliar este fenômeno nas análises
de estabilidade eletromecânica.}
\fonte{EPE-DEE-RE-071/2025-rev0, p.~203.}
\end{citacao}

\begin{destaque}{O detalhe que importa para quem vai implementar}
O primeiro escalonamento do benchmark \textbf{falhou}: o modelo \enquote{mostrou desempenho
insuficiente para resposta a faltas na linha CC, apresentando sobretensões inadequadas e
irrealistas} (p.~199). Só se tornou funcional após a inserção de para-raios CC, da lógica de
abertura do NBS e de uma sequência coordenada de eventos de proteção. Essa é a parte difícil do
problema --- e ela está documentada.
\end{destaque}

\subsection{Enunciado}

O país vai construir o maior elo HVDC-VSC em linha aérea já contratado e, hoje:

\begin{itemize}
  \item não existe modelo VSC-MMC aberto, auditável e neutro em relação a fabricante para estudos
        EMT no Brasil;
  \item não existe benchmark público que permita a um terceiro reproduzir ou contestar os
        resultados publicados --- o utilizado pela EPE está dentro de um software proprietário;
  \item os modelos \enquote{realistas} recomendados são, por construção, caixas-pretas de
        fabricante, com restrição de compartilhamento \enquote{por questão de confidencialidade e
        proteção ao segredo industrial} (p.~22, nota 10);
  \item o avanço do CEPEL ocorre no flanco \textbf{RMS/fasorial}, não no flanco EMT.
\end{itemize}

\subsection{Pergunta de pesquisa}

\begin{destaque}{Pergunta de pesquisa}
É possível construir um modelo genérico de conversor HVDC-VSC (MMC meia-ponte), implementado em
ATP-EMTP e acoplado à rede por fonte controlada por MODELS, publicado sob licença permissiva e
integrado ao Olivas EPS, capaz de: (i)~reproduzir um modelo chaveado detalhado construído no próprio
ATP com NRMSE $\leq$~2\%; (ii)~reproduzir o balanço de regime permanente publicado pela EPE com erro
$\leq$~1\%; e (iii)~reproduzir, \emph{em envelope paramétrico}, os marcos transitórios de falta
polo-terra na linha CC do Capítulo 12 --- tudo isso sem qualquer dado confidencial de fabricante e
com rastreabilidade auditável?
\end{destaque}

%==============================================================================
\section{Objetivos}
%==============================================================================

\subsection{Objetivo geral}

Desenvolver, validar quantitativamente e publicar sob licença permissiva um modelo eletromagnético
transitório genérico de estação conversora HVDC-VSC com submódulos 100\% meia-ponte --- braço com
estados de capacitor observáveis, bloqueio por diodos de braço, controle \emph{grid-following} e
sequência de proteção e religamento de falta CC --- executável em ATP-EMTP e integrado ao Olivas
Power System Studio, acompanhado de um benchmark aberto e reprodutível do bipolo $\pm$600~kV,
3.000~MW e 2.500~km Angicos--Itaporanga~2.

\subsection{Objetivos específicos (partição MoSCoW)}

\paragraph{Obrigatórios (\emph{must}) --- núcleo mínimo defensável.}
\begin{enumerate}[label=\textbf{OE\arabic*.}, leftmargin=1.6cm]
  \item Consolidar as premissas públicas do bipolo em arquivo versionado, com \textbf{campo de fonte
        e faixa de incerteza por parâmetro}, registrando as divergências internas reais do relatório.
  \item Executar, como \textbf{primeira tarefa técnica e com decisão \emph{go/no-go} formal}, a prova
        de conceito do acoplamento MODELS para a rede no ATP, e implementar o acoplamento automático
        que hoje não existe no produto.
  \item Formular e verificar o braço MMC meia-ponte com estados de capacitor por submódulo, com
        oráculo numérico independente em Python/NumPy.
  \item Implementar o estado bloqueado, com condução pelos diodos de braço --- o mecanismo em que o
        modelo escalado da EPE falhou.
  \item Derivar os parâmetros internos não publicados a partir de regras de projeto de literatura
        aberta, com \textbf{faixa declarada}: tensão de submódulo, número de submódulos, constante de
        energia, indutância de braço e impedância do transformador conversor.
  \item Gerar a linha CC de 2.500 km por modelo JMarti, usando o caso de dados LINE CONSTANTS do
        próprio ATP.
  \item Implementar em MODELS o controle \emph{grid-following}: SRF-PLL, laços de corrente em eixo
        \emph{dq}, limitador com base declarada, curva de capacidade P-Q, operação a 85\% da tensão
        nominal e modo STATCOM.
  \item Implementar a sequência de proteção e religamento com \textbf{todas as temporizações
        parametrizadas}.
  \item Publicar repositório aberto com DOI e manifesto SHA256.
\end{enumerate}

\paragraph{Desejáveis (\emph{should}).}
Estudo de agregação de submódulos com fronteira de Pareto entre exatidão e custo computacional;
varredura de sensibilidade das temporizações de proteção; e decomposição do erro de representação do
conversor, comparando o modelo completo à aproximação por carga ZIP 60/40 que a EPE foi obrigada a
adotar.

\paragraph{Opcionais (\emph{could}).}
Integração aditiva ao Olivas EPS, métricas de conversor e codificação executável dos requisitos da
EPE.

%==============================================================================
\section{Delimitação de escopo}
%==============================================================================

\subsection{Incluído}

\begin{itemize}
  \item MMC 100\% meia-ponte, com braço acoplado à rede por \textbf{fonte controlada por MODELS},
        com atraso de um passo declarado e verificado quanto à estabilidade.
  \item \textbf{Modelo chaveado detalhado de referência construído no próprio ATP} (8 a 12
        submódulos por braço, com \emph{snubbers}), usado como verdade-terreno interna --- e
        reexecutável pela banca sem licença de terceiros.
  \item Topologia CC completa: dois polos independentes, reator de alisamento, reator de braço,
        resistor de pré-inserção, \emph{Neutral Bus Switch}, para-raios CC, ramo de aterramento,
        eletrodo de terra e retorno metálico, com os três modos de operação comutáveis.
  \item Linha CC de 2.500 km em JMarti, com verificação de velocidade de modo aéreo e tempo de
        trânsito ($\approx$~8,5~ms).
  \item Transformador conversor via BCTRAN; redes CA por equivalentes de Thévenin com \textbf{SCR
        paramétrico}: 2,50, 4,08, 5,47 e 6,07, com 1,92 reportado.
  \item Varreduras de nível de tensão CC em 525, 600 e 640~kV, cobrindo os três cenários de mercado
        do \S1.5.4 do relatório.
\end{itemize}

\subsection{Excluído, e por quê}

\begin{center}
\small
\begin{tabularx}{\textwidth}{@{}p{5.1cm}X@{}}
\toprule
\textbf{Item excluído} & \textbf{Motivo} \\
\midrule
Componente Type-94 e elementos baseados em compensação &
o modo THEV admite um único elemento por sub-rede eletricamente conectada, e o bipolo tem 24 braços
em, no máximo, duas sub-redes separadas pela linha \\
Controle \emph{grid-forming} &
a própria EPE registra que \enquote{não há uma definição consensada} de requisitos (\S15.2.15,
p.~256); validar contra requisito instável é indefensável --- vira tema-satélite \\
Full-bridge, arranjos mistos, híbridos LCC+VSC, acima de 640~kV, multiterminais &
descartados pelo grupo de trabalho EPE-ONS (\S6.4, p.~79) \\
Disjuntores CC &
não fazem parte da solução meia-ponte de referência \\
Espectro harmônico de comutação e conformidade IEEE 519 &
um modelo com submódulos agregados não reproduz o espectro de comutação; prometer isso destrói a
credibilidade técnica \\
Estabilidade dirigida por conversores, varredura de impedância e passividade &
tema-satélite \\
Coordenação de isolamento e dimensionamento de para-raios, \emph{chopper} e resistores &
tema-satélite; o modelo entrega as formas de onda, não o dimensionamento \\
Modelos \emph{replica} de fabricante e padrão IEEE/CIGRÉ DLL &
tema-satélite e item de roadmap comercial \\
Refatoração do núcleo do Olivas EPS &
a integração é estritamente aditiva \\
Redistribuição do solver ATP &
publicam-se dados e modelos, nunca o executável \\
\bottomrule
\end{tabularx}
\end{center}

%==============================================================================
\section{Metodologia e cronograma}
%==============================================================================

Dois marcos são gates de decisão, e não entregas. O \textbf{M1} decide se a arquitetura é viável
antes de qualquer investimento em controle; o \textbf{M5} é o núcleo mínimo defensável --- sozinho,
já constitui um TCC aprovável e um artigo submetível.

{\small
\begin{xltabular}{\textwidth}{@{}p{2.1cm}p{1.6cm}X@{}}
\toprule
\textbf{Período} & \textbf{Marco} & \textbf{Conteúdo e critério de passagem} \\
\midrule
\endfirsthead
\toprule
\textbf{Período} & \textbf{Marco} & \textbf{Conteúdo e critério de passagem} \\
\midrule
\endhead
\bottomrule
\endfoot
set./2026 & M0 &
Extração rastreável das premissas; benchmark especificado em cinco casos de ensaio; critérios
pré-registrados. \emph{Premissas e requisitos aprovados pelo orientador.} \\
out./2026 & \textbf{M1} \newline \emph{gate} &
\textbf{Go/no-go em 31/10/2026.} Prova de conceito do acoplamento MODELS para fonte controlada;
oráculo numérico em Python; memorial de dimensionamento. \emph{A prova de conceito executa e o laço
realimentado com atraso de um passo é estável.} \\
nov./2026 & M2 &
Linha CC em JMarti; braço em MODELS com estados de submódulo; emissor em formato colunar estrito.
\emph{Erro $\leq$~1\% nos parâmetros de linha; $\tau \approx$~8,5~ms.} \\
dez./2026 & M3 &
Verdade-terreno interna no próprio ATP; teste de conservação de energia de braço; \emph{preprint}
com DOI. \emph{TCC~1 aprovado; NRMSE $\leq$~2\%.} \\
jan./2027 & M4 &
Controle \emph{grid-following}; sintonia por alocação de largura de banda; curva de capacidade;
modo STATCOM. \emph{Degrau de 0,1 a 1,0~pu com sobressinal $\leq$~10\% e acomodação $\leq$~100~ms.} \\
fev./2027 & \textbf{M5} \newline \emph{gate} &
\textbf{Núcleo mínimo defensável, 28/02/2027.} Bipolo completo; inicialização determinística e
\emph{soft-start}; regime permanente validado. \emph{Balanço das Tabelas 15-6 e 15-7 com erro
$\leq$~1\%.} \\
mar./2027 & M6 &
Máquina de estados de proteção; ensaio de falta polo-terra na linha CC. \emph{Data de corte de
escopo. Temporizações reproduzidas com erro $\leq$~2~ms; sobretensão $\leq$~1,5~pu.} \\
abr./2027 & M7 &
Agregação de submódulos; varredura do passo de integração e das temporizações; decomposição do erro
de representação. \emph{Par recomendado de número de submódulos equivalentes e passo; artigo
submetido.} \\
mai.--jun./2027 & M8 e M9 &
Integração aditiva ao Olivas EPS; varreduras de SCR e de tensão CC; redação e defesa.
\emph{Pull request com integração contínua verde; TCC defendido; versão 1.0 com DOI.} \\
\end{xltabular}}

%==============================================================================
\section{Plano de validação}
%==============================================================================

Todos os critérios numéricos são \textbf{pré-registrados} antes da execução dos ensaios, com métrica
única declarada, e \textbf{todos os desvios são publicados}, inclusive os reprovados. Duas classes de
tolerância são mantidas rigorosamente separadas: \emph{verificação de implementação}, que testa o
código e admite tolerância apertada, e \emph{validação física}, que testa o modelo contra grandezas
dependentes de parâmetros internos não publicados e exige tolerância larga com demonstração em
envelope.

\begin{center}
\small
\begin{tabularx}{\textwidth}{@{}p{1.1cm}p{4.3cm}Xp{2.1cm}@{}}
\toprule
\textbf{Nível} & \textbf{O que testa} & \textbf{Critério} & \textbf{Terceiros?} \\
\midrule
V0 & Conservação de energia de braço & desvio $<$~0,5\% da energia nominal em 1~s & não (analítico) \\
V1 & Parâmetros da linha CC & $R$ a 20 e 50~\textcelsius\ com erro $\leq$~1\%; $\tau \approx$~8,5~ms & não \\
V2 & Verdade-terreno construída no próprio ATP & NRMSE $\leq$~2\%; erro de tensão de capacitor $\leq$~1\% & não \\
V3 & Regime permanente (Tabelas 15-6 e 15-7) & erro $\leq$~1\% nas grandezas CC; 2,5~kA e 2,34~kA & não \\
V4 & Desempenho de controle & sobressinal $\leq$~10\%; margem de fase $\geq$~45\textdegree & não \\
V5 & Falta CC (Capítulo 12) & temporizações com erro $\leq$~2~ms; sobretensão $\leq$~1,5~pu; picos \emph{em envelope} & não \\
V6 & Referência externa (Simulink ou PSCAD acadêmico) & complementar, declarado como tal & opcional \\
\bottomrule
\end{tabularx}
\end{center}

\begin{destaque}{A validação não depende de licença de terceiros}
V0 é analítica; V2 usa como referência um modelo construído no próprio ATP; V1, V3 e V5 usam números
publicados pela EPE. V6 é bônus. Isso remove o risco mais comum desse tipo de trabalho: ficar refém
de uma licença que a universidade pode não ter.
\end{destaque}

Uma correção metodológica relevante: \textbf{os parâmetros internos do MMC não são públicos} ---
número de submódulos, capacitância, indutância de braço e dados do transformador. A resposta não é
arbitrar valores e comparar ponto a ponto, e sim \textbf{varrer as faixas de projeto declaradas e
demonstrar que o resultado da EPE cai dentro do envelope produzido}, separando erro de método de erro
de premissa.

%==============================================================================
\section{Integração ao Olivas Power System Studio}
\label{sec:integracao}
%==============================================================================

\subsection{O que já existe (verificado no repositório)}

O Olivas EPS já é um front-end e pré/pós-processador completo de ATP-EMTP:

\begin{itemize}
  \item \arq{app/preprocessor/atp\_format.py} --- emissão em \textbf{formato colunar estrito}
        (até 80 colunas, campo de nó de 6 caracteres), conforme o \emph{ATP Rule Book}, volume 2;
  \item \arq{app/preprocessor/bridge\_to\_atp.py} (2.461 linhas) --- tradutor de unifilar para
        cartões, com despacho por tipo de componente;
  \item \arq{app/preprocessor/spec/} e \textbf{59 arquivos \arq{.ocomp}} em \arq{catalog\_specs/}
        --- registro declarativo de componentes;
  \item \arq{tacs.py}, \arq{tacs\_blocks.py}, \arq{tacs\_tf.py} e \arq{tacs\_interface.py} ---
        blocos TACS e ponte entre o solver e o TACS;
  \item \arq{vcb\_model\_emitter.py} e \arq{atp\_templates/vcb\_reignition.mod} ---
        \textbf{precedente direto}: um algoritmo complexo (CIGRE WG A3.26) já é entregue como
        template em linguagem \textbf{ATP MODELS};
  \item \arq{lcc.py}, \arq{jmarti\_pch.py} e \arq{vector\_fitting.py} --- LINE CONSTANTS, JMarti e
        ajuste racional; \arq{bctran.py} --- transformador;
  \item \arq{simulation/runner.py} e \arq{results\_reader.py} --- execução do ATP e leitura de PL4;
  \item \arq{postprocessor/audit\_trail.py} --- laudo com SHA256 dos dados de entrada e bloco de
        limitações automático.
\end{itemize}

\subsection{O delta do TCC e três achados de auditoria}

Nenhum dos 59 componentes do catálogo é de corrente contínua: não há conversor, não há polo CC, e a
enumeração de fases admite apenas A, B, C, neutro e nenhuma. Além disso:

\begin{center}
\small
\begin{tabularx}{\textwidth}{@{}p{5.4cm}X@{}}
\toprule
\textbf{Achado} & \textbf{Consequência} \\
\midrule
\textbf{Não existe acoplamento MODELS--rede automatizado.} O próprio
\arq{vcb\_reignition.mod} declara que \enquote{o wiring TACS [\ldots] é manual}. &
É o gargalo estrutural e a \textbf{primeira tarefa técnica}. Isoladamente, vale mais para o produto
do que o próprio módulo HVDC: destrava qualquer modelo futuro que precise de lógica --- proteção,
FACTS, religador, gerador com conversor. \\
\arq{tacs\_interface.py} documenta a fonte comandada por TACS como \enquote{SOURCE TYPE-13 (V) ou
TYPE-14 (I)}. &
No ATP, o Type-13 é fonte rampa de dupla inclinação e o Type-14 é fonte cossenoidal; a fonte
comandada por TACS/MODELS é o \textbf{Type-60}. É exatamente o tipo de erro que faz uma prova de
conceito falhar pelo motivo errado. \\
\arq{tests/test\_pp\_registry\_guard.py} espera 49 arquivos \arq{.ocomp}; existem 59. &
Drift pré-existente do teste de guarda, a corrigir em \emph{pull request} separado, e não como
carona do PR de HVDC. \\
\bottomrule
\end{tabularx}
\end{center}

\subsection{Entregáveis de integração, todos aditivos}

Quatro arquivos \arq{.ocomp} (conversor, linha CC, NBS e eletrodo); renderizadores gráficos; ramos de
despacho no tradutor; extensão da enumeração de fases com polos CC e retorno; emissão estrita do
para-raios CC; módulo de métricas de conversor; codificação executável dos requisitos \S3.4[8] e
\S3.4[9] com citação de página; e um estudo em lote com veredito de aprovação encadeado ao laudo
auditável que o produto já emite.

%==============================================================================
\section{Viabilidade de mercado}
\label{sec:mercado}
%==============================================================================

\subsection{A premissa precisa ser corrigida antes de virar argumento}

\textbf{\enquote{Hoje só o PSCAD tem modelo voltado para HVDC-VSC} é falso como afirmação geral.}
Têm modelo MMC nativo, no mínimo: EMTP-RV, com documentação pública de 2014 descrevendo quatro níveis
de modelo; DIgSILENT PowerFactory, com meia-ponte e ponte completa cobrindo os tipos 3 a 7 da CIGRÉ
TB~604; MATLAB/Simulink com Simscape Electrical; PLECS; RTDS/RSCAD; OPAL-RT/HYPERSIM; e Typhoon HIL.
Sustentar a premissa original em banca --- ou em material de venda --- é entregar o flanco ao
primeiro interlocutor que consultar a documentação desses fornecedores.

\medskip
\noindent\textbf{A premissa é verdadeira em três recortes}, e é deles que o trabalho vive:

\begin{enumerate}
  \item \textbf{Ecossistema ATP-EMTP/ATPDraw:} não foi localizado nenhum modelo MMC/VSC público,
        distribuído ou de biblioteca. Há artigos de conferência sem entrega de arquivos.
        \emph{Exigência metodológica: isso só entra no texto como \enquote{não foi localizado até
        [data]}, com script de busca versionado --- tarefa do primeiro mês.}
  \item \textbf{Ecossistema aberto:} Dynawo, DPsim, ANDES, OpenIPSL, PowerSimulationsDynamics.jl,
        GridPACK e pandapower param em modelos fasoriais ou médios. Nenhum MMC EMT chaveado.
  \item \textbf{Prática regulatória brasileira:} nas 509 páginas do relatório da EPE há
        \textbf{15 ocorrências de \enquote{PSCAD}} e \textbf{zero} de ATP, EMTP, ATPDraw,
        PowerFactory, Simulink, RTDS ou HYPERSIM.
\end{enumerate}

\subsection{O argumento forte, verificado na fonte primária}

O gargalo brasileiro \textbf{não é o PSCAD; é o ATP}. Nos Procedimentos de Rede do ONS, Submódulo 2.5
(\emph{Requisitos mínimos para elos de corrente contínua}, revisão 2016.12):

\begin{citacao}
\enquote{Mesmo que os estudos tenham sido realizados por meio do PSCAD, a transmissora deve entregar
ao ONS, até o início dos estudos pré-operacionais, um modelo detalhado do elo CC incluindo todos os
controles também para o programa ATP. Este modelo deve ser acompanhado pelo manual correspondente e
pelos testes de validação executados confrontados com os resultados obtidos pelo Simulador CC.}
\fonte{ONS, Submódulo 2.5, \S7.15.1.10, p.~36.}
\end{citacao}

\begin{citacao}
Sobre os estudos de manobra do pátio CA das conversoras --- energização de transformadores
conversores, tensão de restabelecimento transitória de disjuntores, energização de linhas e
religamentos: \enquote{As simulações devem ser realizadas na ferramenta ATP.}
\fonte{ONS, Submódulo 2.5, \S7.15.3.4, p.~37.}
\end{citacao}

E no Submódulo 18.2 (\emph{Modelos computacionais}, revisão 2020.01), item 4.2.4: \enquote{ATP ---
Denominação de referência: Modelo para análise de transitórios eletromagnéticos. Propriedade: uso
livre}, referenciado em 15 submódulos. \textbf{O PSCAD não consta da lista oficial de ferramentas do
ONS.}

\begin{destaque}{Conclusão}
Existe \textbf{obrigação regulatória} de produzir modelo detalhado de elo CC em ATP e, para o
primeiro bipolo VSC do SIN, alguém terá de construir um MMC em ATP. Hoje esse trabalho é feito
artesanalmente, sem ferramental de autoria, e não há modelo público de partida.

\medskip
\emph{Diligência:} os Procedimentos de Rede passaram por revisões recentes. A numeração acima foi
verificada nas versões citadas; reconfirmar na revisão vigente antes de qualquer uso externo.
\end{destaque}

\subsection{Dimensionamento honesto}

\begin{destaque}{O que não é mercado}
Os R\$~26,5 bilhões de investimento recomendado, dos quais R\$~17,1 bilhões em instalações CC, são
\textbf{CAPEX de obra}. Usá-los como mercado endereçável de software é o erro de dimensionamento mais
fácil de desmontar numa reunião. O gasto brasileiro realmente endereçável em ferramenta de estudo EMT
de HVDC é da ordem de \textbf{poucos milhões de reais por ano, entre menos de dez organizações} ---
a maioria já com PSCAD instalado e, a partir do relatório R2, obrigada a usar modelo \emph{replica}
de fabricante.
\end{destaque}

\noindent\textbf{Onde está o retorno, em ordem de concretude:}

\begin{enumerate}
  \item \textbf{A linha de montagem normativa.} Transformar os requisitos \S3.4[8] e \S3.4[9] em
        critérios executáveis, com veredito de aprovação e laudo auditável. O modelo físico qualquer
        engenheiro competente reproduz da literatura; a codificação normativa versionada e
        rastreável, não.
  \item \textbf{Infraestrutura que vale mais que o módulo.} O acoplamento MODELS--rede destrava
        qualquer modelo futuro que precise de lógica e habilita módulos de terceiros.
  \item \textbf{Alavanca sobre o que já vende:} tensão de restabelecimento transitória segundo a
        IEC 62271-100 nos disjuntores CA das conversoras --- que são precisamente o meio de
        eliminação de falta CC na solução meia-ponte --- e sobretensões agora também do lado CC.
  \item \textbf{Funil acadêmico:} universidades sem licença de PSCAD passam a poder ensinar HVDC-VSC
        com ATP e Olivas EPS. É aquisição barata, não receita.
\end{enumerate}

\textbf{Teto declarado, por escrito e por código:} nenhum modelo genérico aberto substitui o
\emph{replica} do fabricante em relatório R2, detalhamento ou comissionamento. O nicho é
planejamento, anteprojeto, verificação independente e ensino --- exatamente onde a EPE ficou sem
ferramenta. O laudo emite esse bloco de limitações automaticamente.

\textbf{Gate comercial antes de maio de 2027:} lista nominal de dez organizações-alvo, hipótese de
preço por assento ou por estudo, e duas cartas de interesse assinadas. Sem isso, o módulo deve ser
orçado como marketing técnico e funil acadêmico --- o que continua sendo bom investimento, com a
expectativa correta.

\begin{destaque}{Um risco jurídico que precede este TCC}
A licença de uso do ATP-EMTP é concedida individualmente por grupos regionais de usuários e exige
declaração de \textbf{não participação no \enquote{comércio do EMTP}}, além de vedar a
redistribuição do material. Isso é exposição \textbf{pré-existente e de nível de empresa}: a Olivas
já comercializa hoje um front-end de ATP, sem qualquer módulo HVDC. O parecer jurídico e a consulta
ao grupo de usuários devem ser conduzidos pela empresa agora, por causa do produto existente, e
nenhuma peça deste trabalho deve alegar que o módulo HVDC agrava a exposição. O que o TCC publica são
\textbf{dados e modelos}, nunca o solver.
\end{destaque}

%==============================================================================
\section{Estratégia de licenciamento}
%==============================================================================

\begin{center}
\small
\begin{tabularx}{\textwidth}{@{}p{2.9cm}Xp{2.6cm}@{}}
\toprule
\textbf{Camada} & \textbf{Conteúdo} & \textbf{Licença} \\
\midrule
\textbf{Aberta} \newline (repositório com DOI) &
Modelo, templates MODELS, cartões ATP, oráculo numérico, premissas, conjunto de dados de validação,
figuras e documentação &
Apache-2.0 \newline CC-BY-4.0 \\
\textbf{Produto} \newline (Olivas EPS) &
Estudo automatizado com veredito de aprovação, codificação normativa, motor de varredura
paramétrica, laudo profissional com auditoria e SHA256, biblioteca curada de casos e suporte &
proprietária \\
\textbf{Infraestrutura} &
Acoplamento MODELS--rede e registro de emissores, contribuídos ao núcleo de forma estritamente
aditiva &
Apache-2.0 \\
\bottomrule
\end{tabularx}
\end{center}

Racional: a barreira de entrada de um concorrente \textbf{não é o modelo} --- é a cadeia
\emph{unifilar $\rightarrow$ cartão ATP $\rightarrow$ execução $\rightarrow$ laudo assinável}. Abrir o
modelo maximiza adoção, citação e confiança; o produto captura valor no fluxo de trabalho.

%==============================================================================
\section{Riscos e mitigação}
%==============================================================================

{\footnotesize
\begin{xltabular}{\textwidth}{@{}p{4.4cm}p{1.1cm}p{1.9cm}X@{}}
\toprule
\textbf{Risco} & \textbf{Prob.} & \textbf{Impacto} & \textbf{Mitigação} \\
\midrule
\endfirsthead
\toprule
\textbf{Risco} & \textbf{Prob.} & \textbf{Impacto} & \textbf{Mitigação} \\
\midrule
\endhead
\bottomrule
\endfoot
Acoplamento MODELS para fonte controlada não funcionar, ou o laço com atraso de um passo ser instável
& média & alto &
Prova de conceito como primeira tarefa, com \emph{go/no-go} em 31/10/2026, antes de qualquer
investimento em controle \\
Custo computacional inviabilizar as varreduras, já que MODELS é interpretado
& média & médio &
Envelope declarado desde o início: submódulos agregados, passo de 20~$\mu$s e reordenação apenas nas
transições de nível \\
Nenhum parâmetro interno do MMC é publicado
& certa & médio &
Dimensionamento por regras públicas com faixa declarada, e validação demonstrada em envelope \\
Modelo não reproduzir a falta CC --- falha que o próprio benchmark do PSCAD apresentou
& média & alto &
Estado bloqueado validado isoladamente em circuito reduzido antes do caso completo \\
Escopo excessivo para autor único em dez meses
& alta & alto &
Partição MoSCoW assinada e gate M5, que sozinho já constitui TCC aprovável \\
Perda de ineditismo durante a execução
& média & médio &
Revisão exaustiva do estado da arte no primeiro mês; \emph{preprint} com DOI em dezembro de 2026
para fixar prioridade \\
Indisponibilidade de ferramenta de referência externa
& média & baixo &
A validação principal é autossuficiente (V0 a V5) \\
Cláusula de comércio do EMTP na licença do ATP
& média & alto para o negócio; nulo para a validade acadêmica &
Parecer jurídico conduzido pela empresa agora, por causa do produto existente \\
Modelo aberto ser percebido como substituto do \emph{replica} de fabricante
& média & médio &
Teto declarado por escrito \textbf{e por código}: bloco de limitações automático no laudo \\
Quebra da integração contínua na integração ao produto
& alta & baixo &
Correção do drift do teste de guarda em PR separado; código novo importável sem dependência de GUI \\
\end{xltabular}}

%==============================================================================
\section{Questões deliberadamente em aberto}
%==============================================================================

Levantadas na revisão crítica desta proposta e \textbf{não resolvidas aqui} --- são o conteúdo das
primeiras seis semanas de trabalho.

\begin{enumerate}
  \item \textbf{Taxonomia do modelo.} Um braço com vetor de tensões de capacitor, modulação por nível
        mais próximo e balanceamento por ordenação é mais próximo do \emph{Detailed Equivalent Model}
        \cite{gnanarathna2011} com submódulos agregados do que do \emph{Switching Function of Arm}.
        A rotulagem correta muda a alegação de contribuição e precisa ser fechada antes da redação.
  \item \textbf{Topologia do estado bloqueado.} Um diodo de \emph{bypass} em paralelo com uma fonte de
        tensão ideal controlada curto-circuita a fonte ao conduzir, produzindo matriz singular. A
        implementação viável provavelmente é chavear o \emph{valor} da fonte dentro do MODELS conforme
        o sinal da corrente de braço, com resistência série explícita. É onde o trabalho pode queimar
        semanas --- e onde o modelo da EPE falhou.
  \item \textbf{Tipo de cartão da fonte comandada:} Type-60, e não Type-13 ou Type-14 como o
        repositório documenta.
  \item \textbf{Requisito de reenergização.} A leitura de que \enquote{a EPE não atenderia ao próprio
        requisito} \textbf{não se sustenta}: o \S15.2.13 (p.~256) mostra que os 2,5~s e os 5~s foram
        \emph{derivados} dos aproximadamente 2~s de reenergização e 4~s de recuperação observados no
        PSCAD. Resta uma ambiguidade de definição de \enquote{reenergização} --- nota de rodapé, e não
        pilar de originalidade.
  \item \textbf{Divergências internas do relatório} a registrar, e não a \enquote{explicar}: 3~kA
        (p.~254) contra aproximadamente 2,85~kA (p.~123) de corrente máxima de IGBT; \S3.4[8].b
        (\enquote{corrente total} não superior a 1~pu) contra \S15.2.9 (limite de 1~pu de corrente
        \emph{CC}); condutor Kiwi com formação 72/7 contra 45/7 entre tabelas; e
        \enquote{R\$~199,96/MWh} (p.~242) contra \enquote{/MW} (p.~246 e 259).
\end{enumerate}

%==============================================================================
\section{Temas-satélite}
%==============================================================================

Próximos trabalhos sobre a mesma base. O modelo aberto na versão 1.0 é pré-requisito de todos.

\begin{center}
\small
\begin{tabularx}{\textwidth}{@{}p{4.6cm}Xp{2.7cm}@{}}
\toprule
\textbf{Tema} & \textbf{O que ataca} & \textbf{Perfil} \\
\midrule
\textbf{Grid-forming sobre o modelo aberto} --- droop, máquina virtual síncrona e dVOC em MODELS &
O regime em que o \emph{grid-following} falha por construção: SCR abaixo de 2, exatamente a
contingência N-2 Açu~III--Angicos &
Graduação com base forte em controle, ou mestrado \\
\textbf{Estabilidade dirigida por conversores em Angicos} --- varredura de impedância e passividade &
Recomendação \S3.4[11] e \S15.2.14: interações de controle e harmônicas com os parques eólicos e
solares vizinhos &
Mestrado; maior valor comercial \\
\textbf{Curto-circuito de unidade com conversor pleno} --- IEC 60909-0:2016, \S6.7 &
Fecha o ciclo com o módulo de curto-circuito que o Olivas EPS já comercializa e que hoje não
implementa o \S6.7 &
Graduação; maior conversão imediata \\
\textbf{Coordenação de isolamento do pátio CC $\pm$600~kV} &
Recomendação \S3.4[10]: para-raios, resistor de amortecimento e \emph{chopper} &
Graduação com perfil de alta tensão \\
\textbf{Padrão IEEE/CIGRÉ DLL} (TB~958, 2025) em fluxo ATP-EMTP &
O único caminho que \emph{resolve}, em vez de substituir, o problema do modelo certificado de
fabricante &
Roadmap de produto, não pesquisa de graduação \\
\bottomrule
\end{tabularx}
\end{center}

%==============================================================================
\section{Decisões pendentes}
%==============================================================================

\begin{enumerate}
  \item \textbf{Orientador e instituição} --- o trabalho exige alguém confortável com transitórios
        eletromagnéticos e com o ATP.
  \item \textbf{Ferramenta de referência externa} --- verificar se a universidade dispõe de PSCAD
        educacional ou de MATLAB com Simscape. Não é bloqueante, mas fortalece a validação.
  \item \textbf{Declaração de vínculo} --- o autor é ligado à Olivas; isso deve constar da folha de
        rosto e de seção própria da monografia.
  \item \textbf{Parecer jurídico sobre a licença do ATP} --- tarefa da empresa, imediata e
        independente do TCC.
  \item \textbf{Onde publicar o repositório aberto} --- organização própria ou institucional, com DOI
        no Zenodo.
\end{enumerate}

%==============================================================================
\begin{thebibliography}{99}
\addcontentsline{toc}{section}{Referências}

\bibitem{epe2025}
EMPRESA DE PESQUISA ENERGÉTICA. \textbf{EPE-DEE-RE-071/2025-rev0 --- Estudo de Expansão das
Interligações Regionais, Parte III: Expansão da Capacidade de Exportação da Região Nordeste e
Importação da Região Sul}. Rio de Janeiro: EPE, nov. 2025. 509~p.

\bibitem{ons25}
OPERADOR NACIONAL DO SISTEMA ELÉTRICO. \textbf{Procedimentos de Rede --- Submódulo 2.5: Requisitos
mínimos para elos de corrente contínua}. Revisão 2016.12.

\bibitem{ons182}
OPERADOR NACIONAL DO SISTEMA ELÉTRICO. \textbf{Procedimentos de Rede --- Submódulo 18.2: Modelos
computacionais}. Revisão 2020.01.

\bibitem{epe2020}
EMPRESA DE PESQUISA ENERGÉTICA. \textbf{EPE-DEE-DEA-NT-004/2020 --- Diretrizes para a elaboração dos
relatórios técnicos para a licitação de novas instalações da Rede Básica: estrutura e conteúdo dos
relatórios R1, R2, R3, R4 e R5}. 2020.

\bibitem{gnanarathna2011}
GNANARATHNA, U.~N.; GOLE, A.~M.; JAYASINGHE, R.~P. Efficient Modeling of Modular Multilevel HVDC
Converters (MMC) on Electromagnetic Transient Simulation Programs. \textbf{IEEE Transactions on Power
Delivery}, v.~26, n.~1, p.~316--324, 2011.

\bibitem{cigre604}
CIGRÉ WORKING GROUP B4.57. \textbf{Technical Brochure 604 --- Guide for the Development of Models for
HVDC Converters in a HVDC Grid}. Paris: CIGRÉ, 2014.

\bibitem{cigre958}
CIGRÉ JOINT WORKING GROUP B4.82. \textbf{Technical Brochure 958 --- Guidelines for Use of Real-Code in
EMT Models for HVDC, FACTS and Inverter-Based Generators in Power Systems Analysis}. Paris: CIGRÉ,
2025.

\bibitem{saad2014}
SAAD, H.; MAHSEREDJIAN, J.; DENNETIÈRE, S. et al. \textbf{MMC Documentation}. EMTP-RV, 2014.

\bibitem{ieee2021}
HATZIARGYRIOU, N. et al. Definition and Classification of Power System Stability --- Revisited and
Extended. \textbf{IEEE Transactions on Power Systems}, v.~36, n.~4, p.~3271--3281, 2021.

\bibitem{dommel}
DOMMEL, H.~W. \textbf{EMTP Theory Book}. Portland: Bonneville Power Administration.

\bibitem{rulebook}
LEUVEN EMTP CENTER. \textbf{ATP Rule Book}, volumes 1 e 2.

\end{thebibliography}

\end{document}
